\newcommand{\gibbs}{P_t^{\b,\om}}
\newcommand{\gibbss}{P_s^{\b,\om}}
\newcommand{\gibbsDeux}{P_t^{\b,\om,\otimes 2}}

% SF stuff
\newcommand{\gibbsni}{P_{n-i}^{\b,\theta_{i,x}\om}}
\newcommand{\gibbsmn}{P_{m+n}^{\b,\om}}
\newcommand{\gibbsmdn}{P_{m}^{\b,\theta_{n,y}\om}}
\newcommand{\gibbsdn}{P_{2n}^{\b,\om}}
\newcommand{\gibbsbu}{P_n^{\b_1,\om}}
\newcommand{\gibbsbd}{P_n^{\b_2,\om}}
\newcommand{\gibbsbp}{P_n^{\b',\om}}
\newcommand{\gibbsdnpu}{P_{2n+1}^{\b,\om}}
\newcommand{\gibbsmu}{P_{t-1}^{\b,\om}}
\newcommand{\gibbsc}{P_{t}^{\b,\om,2}}
\newcommand{\gibbsmuc}{P_{t-1}^{\b,\om, 2}}
\newcommand{\gibbspc}{[P_{t}^{\b,\om}]^{\otimes 2}}
\newcommand{\gibbsmupc}{[P_{t-1}^{\b,\om}]^{\otimes 2}}
\newcommand{\gibbsp}{P_n^{\b',\om}}
% end SF stuff





 \chapter{Free energy and phase transition} \label{sec:thermo}
 
 

 
 %%%%%%%%%%
 
 
 Notations and conventions: all through the notes, we will use the same symbols $P, \P, \dots$ to denote probability measures and mathematical expectations; e.g., $P[X]$ is the 
 $P$-expectation of the random variable $X$.  
\medskip

In this section, we introduce the model and two central thermodynamic quantities, the quenched and the annealed free energies.


\section{Polymer model}

The model is defined as a Brownian motion in a random potential. 
\medskip


$\bullet$ {\it The free measure :}   
$(B=\{ B_t\}_{t \geq 0}, P_x)$ is a Brownian motion on 
the $d$-dimensional Euclidean space $\rd$ starting from $x \in \rd$.
We will use short notation  $P_0=P$.

$\bullet$ {\it The random environment:}  
$\om = \sum_i \delta_{(T_i,X_i)}$ is a Poisson point process on $\R_+ \times \rd$ with intensity measure $\nu dt dx$, where $\nu$ is a positive parameter. We suppose that $\om$ is defined on some probability space $(\Omega,\mathcal{G},\IP)$, and we define $\mathcal{G}_t$ to be the $\sigma$-field generated by the environment up to time $t$:
\begin{equation}
\om_t=\om_{|(0,t]\times \rd} \;,\quad \mathcal{G}_t = \sigma\left(\om_t(A) ; A \in \mathcal{B}(\mathbb{R}_+ \times \mathbb{R}^d) \right),
\end{equation}
where $\mathcal{B}(\mathbb{R}_+ \times \mathbb{R}^d)$ denotes the Borel sets of $\mathbb{R}_+ \times \mathbb{R}^d$.
\medskip
%%%%%%%%


From these two basic ingredients, we define the object we consider in the notes. Fix $r>0$, and let $U(x)$ denote Euclidean (closed) ball in $\rd$ with radius $\gamma_d^{-1/d} r$,
$$
U(x)= {\mathrm B}  (x , \gamma_d^{-1/d} r).
$$ with $\gamma_d$ the volume of the unit ball, so $U(x)$ has volume $r^d$.  The tube around path $B$ is the following subset of $(0,t]\times\R^d$: 
\begin{equation}
V_t(B)= \left\{ (s,x): s \in (0,t], x \in U(B_s)\right\}\,.
\end{equation}
When the indicator function
 \begin{equation} \label{def:chitx}
\chi_{s,x}= {\mathbf 1} \{x \in U(B_s) \} = {\mathbf 1} \{|x-B_s| \leq \gamma_d^{-1/d} r\}
\end{equation}
has value 1 [resp., 0],  the path $B$ does see [resp. does not see] the point $(s,x)$. For a fixed path $B$, the quantity defined by
\begin{equation} \label{eq:214}
\om (V_t) = \int_{(0,t] \times \rd} \chi_{s,x} \, \om(ds,dx),
\end{equation}
is the number of Poisson points seen by the path $B$ up to time $t$, playing the role of $\#_t$ in the Introduction. Note that under $\P$, the variable $\om (V_t)$ is Poisson distributed with mean $\nu t r^d$.


$\bullet$ {\it The polymer measure:} Fixing a realization $\om$ of the Poisson point process and a value of the time horizon
$t>0$, we define the probability measure $\gibbs$ on the path space 
${\cC}(\R_+; \rd)$ equipped with its Borel field by  
\begin{equation} \label{mnen}
%%%\m_n (d\w )=P[\z_n: d\w]/P[\z_n],
d \gibbs =\frac{1}{Z_{t}(\om,\b,r)}
\exp \{\beta \om(V_t)\} \;dP,
\end{equation}
where $\b \in \R$ is a parameter (the inverse temperature), where 
\begin{equation} \label{Zt}
Z_t=Z_{t}(\om,\b,r)
=P\lef[\exp \lef(\b \om(V_t)
\rig) \rig]
\end{equation}
is the normalizing constant making $\gibbs$ a probability measure on the path space.
%

 The model has been introduced by Nobuo Yoshida as a polymer model, and first appeared in \cite{CY05} in the literature.  For $\b >0$ the path is attracted by the Poisson points, and repelled otherwise. The Poisson environment represents randomly dispatched impurities. 
For negative $\b$ the model relates to Brownian motion in Poissonian obstacles
\cite{donskervaradhan-wienersaus, Sznitman} which can be traced back to works of Smoluchowski \cite{Smoluchowski}. Here we consider  a directed version, in contrast to crossings \cite{Wuthrich1, Wuthrich2, Wuthrich-geodesics, Wuthrich-bornes} where the path is stretched ballistically.
Our model with $\b \to +\8$ is related to Euclidean first passage percolation \cite{HowardNewman97, Howard00} with exponent $\alpha=2$ therein.

Also, for a  branching Brownian motion in random medium \cite{Shiozawa-clt, Shiozawa-loc}, $Z_t$ is equal to the mean population size in the medium given by $\omega$.

\section{Some key formulas and notations}
We first recall three basic formulas that we will use repeatedly.
\begin{itemize}
\item 
For all non-negative and all non-positive measurable functions $h$ on $\mathbb{R}_+\times \mathbb{R}^d$, the Poisson formula for exponential moments (chapter 3. of \cite{LastPenrose}) writes
\begin{equation} \label{eq:PoissonExpMoment}
\P\left[ e^{\int h(s,x) \om_t(ds dx)}\right] = \exp \int_{]0,t]\times \rd} \nu ds dx \left( e^{h(s,x)}-1 \right)   \;.
\end{equation}
%The formula remains true when $h$ is replaced by $i h$, for any real measurable function $f$. \bfc 
The formula remains true when $h$ is replaced by $i h$, for any real integrable function $h$. %\ec
\item Introducing the notation 
\begin{equation}
\lambda(\b)= e^\b-1\;,
\end{equation}

the linearization formula for Bernoulli writes
\begin{equation}\label{eq:linBer}
e^{\b {\bf 1}_A}-1 = (e^\b-1) {\bf 1}_A =  \lambda(\b) {\bf 1}_A\;.
\end{equation}
\item For all $s\geq 0$, we have 
\begin{equation}\int_{\mathbb{R}^d} \chi_{s,x} \hspace{0.3mm} \rmd x = r^d.
\end{equation}
\end{itemize}

\section{Quenched free energy}
It is defined as the rate of growth of the partition function, and it is a self-averaging property.

\begin{theorem} \label{th:freeenergy}
The quenched free energy
$$ p(\b, \nu) =  \lim_{t \to \8} \frac{1}{t} \ln Z_t (\om, \b, r)
$$
exists a.s. and in $L^p$-norm for all $p \geq 1$, and is deterministic,
$$
 p(\b, \nu) = \sup_{t >0}  \frac{1}{t}\P[\ln Z_t]\;.
$$
\end{theorem}
%%
\begin{remark} We omit the parameter $r>0$  from the notation for the free energy.
The reason is that,  in contrast to $\b$ and $ \nu$, it is kept fixed most of the time.
\end{remark}
%%

$\Box$ Let $\theta_{t,x}$ the space-time shift operator on the environment space, 
$$\theta_{t,x} \big(\sum_i \delta_{(T_i,X_i)}\big)= \sum_i \delta_{(T_i-t,X_i-x)}\;.$$ By Markov property of the Brownian motion,  we have for $s, t \geq 0$,
\begin{eqnarray} \nn
{Z_{t+s}} &=& P\left[ e^{\b \om(V_t)} e^{\b \om(V_{t+s} \setminus V_t)} \right] \\ \nn
&=& P\left[ e^{\b \om(V_t)} P\left[ e^{\b \om(V_{t+s} \setminus V_t)} \big \vert B_t\right] \right] \\ \label{eq:Markov2}
&=& P\left[ e^{\b \om(V_t)}  Z_s \circ \theta_{t,B_t}
\right] \\ \label{eq:Markov}
&=& {Z_t} \times \gibbs[ Z_s \circ \theta_{t,B_t}] \;,
\end{eqnarray}
a remarkable identity expressing the Markov structure of the model.
Let $u(t)= \P[ \ln Z_t]$. By the independence property of Poisson points, $\om_{|]s,t]}$ is independent of $\cG_s$ for all $0\leq s \leq t$.  Then, denoting by $\P^{\cG_t}$ the conditional expectation and  conditional
probability given $\cG_t$, we have 
%using independence  properties of the Poisson process,
\begin{eqnarray*}
u(t+s) %&=& \P[  \ln \frac{Z_{t+s}}{Z_t} ]+ \P \ln Z_t \\
&=&  \P \big[ \ln \gibbs[ Z_s \circ \theta_{t,B_t)}]  \big]+ \P \ln Z_t 
\\
&\stackrel{\rm Jensen}{\geq} 
& \P  \gibbs[   \ln  Z_s \circ \theta_{t,B_t)} ]+u(t)  \\
%&\geq &  \P [ \ln  \P^{\cG_t} \gibbs[ Z_s \circ \theta_{t,B_t)}] ]+u(t)  \qquad \rm{(conditional\; Jensen\; ineq.)}\\
&=&  
 \P \P^{\cG_t} \gibbs[   \ln  Z_s \circ \theta_{t,B_t)} ]+u(t)  \\
&\stackrel{\rm Fubini}{=}&  \P  \big[  \gibbs[ \P^{\cG_t}[\ln Z_s \circ \theta_{t,B_t)}] ] \big]+u(t)\\
%&=&  \P [ \ln   \gibbs[ \P[Z_s \circ \theta_{t,B_t)}] ]]+u(t) \\
&=& \P \big[   \gibbs[ u(s) ]\big]+u(t) \qquad\qquad \qquad (\om {\rm \ shift\ invariant})\\
%&=&  \ln   \P[Z_s] +u(t) \\
&=& u(s)+u(t)\;.
\end{eqnarray*}
Hence the function $u(t)$ is superadditive. By the superadditive lemma, we get the existence of the limit
\begin{equation}
\nn
\lim_{ t \to \8} \frac{u(t)}{t} = \sup_{t >0} \frac{u(t)}{t} \;.
\end{equation}
Now, anticipating the concentration inequality \eqref{eq:concentrationlogZt} and the continuous time bridging
\eqref{eq:OrderOfCenteredlogZt},
% cor 5.1.3
we derive that 
\begin{equation}
\nn
\frac{1}{t} \big( \ln Z_t- \P[\ln Z_t]\big) \longrightarrow 0
\end{equation}
almost surely and in $L^p$ for all $p\geq 1$ finite.  \qed
%%
\section{Annealed free energy and hierarchy of moments}
We compute the expectation of the partition function over the medium using \eqref{eq:214} and Fubini,
\begin{eqnarray}
\P[ Z_t] &=&  P \P\left[ e^{\b \int \chi_{s,x} \om_t(ds,dx)} \right] \nn \qquad \qquad \qquad \qquad
%{\rm (Fubini)} 
\\&\stackrel{\eqref{eq:PoissonExpMoment}}{=}&
P \exp \int_{]0,t]\times \rd} \left( e^{\b \int \chi_{s,x}}-1 \right) \nu ds dx \nn \qquad 
%({\rm by\; } \eqref{eq:PoissonExpMoment})
 \\
&\stackrel{\eqref{eq:linBer}}{=}&P \left[  \exp \la(\b)  \int_{]0,t]\times \rd}  \chi_{s,x}  \nu ds dx \right] \nn  \qquad  \qquad 
%({\rm by\; }\eqref{eq:linBer}) 
 \\&= &\label{eq:annealed}
 \exp\{ t \nu \la r^d\} \;.
\end{eqnarray}
%%%%
Hence $\P[ Z_t]$ grows in time at exponential rate  $p^{(1)}(\b, \nu)= \nu \lambda(\b)r^d$. More generally, it is natural to consider the rate of growth of the $s$-th moment 
of the partition function,
$$
p^{(s)}(\b, \nu)= \lim_{t \to \8} \frac{1}{st} \ln \P[ Z_t^s]\,,\quad s >0.
$$
By H\"older inequality, $\|Z\|_r \leq \|Z\|_s$ for $r \leq s$, these rates are non-decreasing in $s$, and for integer values, they can be expressed by handy variational formulas
using large deviation theory. By Jensen's inequality, we have for all $t$
\begin{equation} \label{eq: annealedbound}
\frac{1}{t}\P[\ln Z_t] \leq   \frac{1}{t} \;\ln \P[ Z_t] = p^{(1)}(\nu,\b)\;,
\end{equation}
yielding the so-called {\em annealed bound} : $$ p(\b, \nu) \leq  p^{(1)}(\nu,\b)\;.$$  Summarizing the above,  we have a chain of inequalities
\begin{equation}
\nn
  p(\b, \nu) \leq  p^{(1)}(\b, \nu) \leq \ldots \leq  p^{(k)}(\b, \nu) \leq  p^{(k+1)}(\b, \nu) \leq \ldots . 
\end{equation}
It is commun folklore that in a large class of models, the first inequalities in the above chain are equalities, while they become strict from
$k^*= \inf \{k \geq 0: p^{(k)}(\b, \nu) <  p^{(k+1)}(\b, \nu)$ (with the convention $p^{(0)}=p$). Considering the sequence of rates $(p^{(k)}; k \geq 1)$ is
classical approach to intermittency \cite{CarmonaMolchanov, khoshnevisan, Zeldovich} and sect. 2.4 of \cite{BertiniCancrini95}. %Analysis of Stochastic Partial Differential Equations, Koshnevisan,AMS 2014}.

 In the directed case, we focus at $k=0,1$ only, since the latter is explicit.
%
%%%%%%%%%%%%%%%%%%%%%%%%%%%%%
\begin{proposition} \label{prop:p}
%$\phantom{}$
Basic properties of the free energy:
\begin{enumerate}
\item \label{eq:p1} For $\b \neq 0, \nu >0$, we have
$ \b \nu r^d < p(\b, \nu) \leq \nu \lambda(\b)r^d$.
\item \label{eq:p2}
$ \b \to p(\b, \nu)$ is convex.
\item \label{eq:p3} The excess free energy 
\begin{equation}
\label{def:excessfreeenergy}   
 \psi(\b,\nu)=\nu \lambda(\b)r^d- p(\b, \nu)
\end{equation}
is non-decreasing in $|\b|$ and in $\nu$. It is jointly continuous. 
\end{enumerate}
\end{proposition}
%%%%
\noindent
$\Box$ The second inequality in item \ref{eq:p1} is the annealed bound. The first one follows from an inifinite-dimensional version of Jensen's
inequality; this version being curiously overlooked in the literature, we recall the full statement:
%%%%
\begin{lemma}[Lemma A.1 in \cite{RASaihp11}]\label{lem:RAS-J}
Let $g$ be a bounded measurable function on a product space ${\cal X} \times  {\cal Y}$, $\mu$ a probability measure on $\cal X$ and $\rho$ a probability measure on $\cal Y$. Then
$$
\ln \int_{\cal X} e^{\int_{\cal Y}  g(x,y) d\rho (y) } d \mu(x) \leq \int_{\cal Y}  \left[ \ln \int_{\cal X}  e^{g(x,y)} d\mu (x) \right] d \rho(y) .
$$
\end{lemma}
We apply it with $\rho=\P, \mu=P, g(x,y)=\b \om(V_t)$ to get the desired bound\footnote{We explain in this note 
why the Lemma is an infinite-dimensional version of Jensen's inequality:  
the functional $\psi(f)= \ln \int_{\cal X} e^{f(x)} d\mu(x)$ is convex, and the function $f(\cdot)=g(\cdot, y)$ is randomly chosen with $\rho(dy)$.}.
However this bound is not so great here, since the simple one $p(\b,\nu) \geq t^{-1} \P[ \ln Z_t]$ for a fixed $t$ (which comes from superadditivity of $u(t)$)
is not linear, but strictly convex in $\b$ and then already better.
%
%
%$$
%P\left[\exp (\b \om(V_t)) \right] \geq \exp \left(\b P[\om(V_t)]\right) 
%$$
%with strict inequality if $\b \nu \neq 0$, and in that case,
%$$
% p(\b, \nu) \geq \frac{1}{t}\P[\ln Z_t]  > \frac{\b}{t} \P  P[\om(V_t)]= \b \nu r^d.
%$$   

Item \ref{eq:p2} is the standard convexity of free energy, 
$$
\frac{\partial^2}{\partial \b^2} \ln Z_t = {\rm Var}_{\gibbs} \big( \om(V_t)) \big) > 0\;,
$$
where ${\rm Var}_{\gibbs}$ denotes the variance under the polymer measure in a fixed environment $\om$.

We now turn towards item \ref{eq:p3}, in the case $\b \geq 0$ (the other case being similar). We use specific properties of the medium, infinite divisibility: 
for $\nu, \D>0$, we note that the superposition 
$\om+\hat \om$ of two independent PPP with intensities $\nu$ and $\D$ is  a PPP with intensity $\nu+\D$. Writing $\P$ the expectation 
over both variables ${\om,\hat \om}$, we compute by conditioning
%$\P_\om$ for integration over $\om$  alone and $\P_{\om,\hat \om}$ for integration over both variables, we have
\begin{eqnarray} \nn
\P \ln Z_t(\om) & \stackrel{\b \geq 0}{\leq} & 
\P \ln Z_t(\om + \hat \om) \\ &=& \P \P\left[ \ln Z_t(\om + \hat \om) \vert \om\right] \nn \\
&\stackrel{\rm Jensen}{\leq} & \P \ln \P\left[ Z_t(\om + \hat \om)  \vert \om\right] \nn \\
&=& \P \ln Z_t (\om) + t \D  \lambda(\b) r^d. \nn
\end{eqnarray}
This proves  monotonicity of $\psi$  in $\nu$. This proves  at the same time  continuity in $\nu$ (locally uniformly in $\b$) and the joint continuity in $(\b, \nu)$.
% \bfc [FAUT-IL a ce stade DONNER la formule de densite du PPP($\nu$) par rapport a PPP(1) pour faire la preuve ???] \ec

\medskip


The plain identity $\P [Yf(Y)]= \theta \P[f(Y+1)]$ for a r.v. $Y$ distributed as a Poisson law with mean $\theta$ has a counterpart for PPP,
an integration by parts formula known as Slivnyak-Mecke formula (e.g., p.50 in \cite{KendallStoyanMecke} or th.~4.1~in \cite{LastPenrose}):  Let $\cM$ be the space of point measures on $(0,t] \times \rd$ and $h: 
\R_+ \times \rd \times {\cM} \to \R_+$ measurable, then
\begin{equation}
\label{eq:IPP}
\P \left[ \int h(s,x; \om_t) \om_t(ds,dx) \right]= \int_{(0,t]\times \rd}  \P\big[ h(s,x; \om_t + \delta_{s,x}) \big]\; \nu ds dx  \;.
\end{equation}
With this in hand, we can show monotonicity of $\psi$ in $\b$: 
\begin{eqnarray} 
\frac{\partial }{\partial \b} \P \big[ \ln Z_t \big]
 \nn &=& \P \gibbs [\om(V_t)]   \\
 \nn &=& \P \int \om_t(ds dx)  \frac{ P\left[ \chi_{s,x} e^{\b \om(V_t)} \right]}{Z_t}\\
\nn &\stackrel{\eqref{eq:IPP}}{=}& \P \int_{(0,t]\times \rd} \nu ds dx  \frac{ P\left[ \chi_{s,x} e^{\b (\om(V_t)+\delta_{s,x})} \right]}{P\left[ 
e^{\b (\om(V_t)+\delta_{s,s})} \right]}  \\
\nn &\stackrel{\eqref{eq:linBer}}{=}& \P \int_{(0,t]\times \rd} \nu ds dx  \frac{ P\left[ e^\b \chi_{s,x} e^{\b \om(V_t)} \right]}{P\left[ \big(\lambda(\b)\chi_{s,x}
+1\big) e^{\b \om(V_t)} \right]} \\
\label{eq:meanenergy} 
&{=}& \P \int_{(0,t]\times \rd} \nu ds dx  \frac{ e^\b \gibbs[\chi_{s,x}]}{1+ \lambda(\b) \gibbs[\chi_{s,x}]} \;.
\end{eqnarray}
Define 
\begin{equation}
\label{def:psit}
\psi_t(\b,\nu) = t^{-1} \P \left[ \nu \lambda r^d -\ln Z_t \right] .
\end{equation}
With the identity
$$ \frac{\partial }{\partial \b} \nu \lambda(\b)r^dt = \nu e^{\b} r^d t =  e^{\b} \int_{(0,t]\times \rd} \nu ds dx \gibbs[ \chi_{s,x}]$$
we obtain
\begin{equation}
\label{eq:psi'}
\frac{\partial }{\partial \b}  \psi_t(\b,\nu) = \frac{1}{t}e^\b \lambda(\b) \nu 
 \int_{(0,t]\times \rd}  ds dx \; \P \frac{ \gibbs[  \chi_{s,x}]^2}{1+ \lambda(\b) \gibbs[\chi_{s,x}]} \;,
 \end{equation}
which has the sign of $\b$. So the limit $\psi$ of $\psi_t$ is increasing in $|\b |$. \qed
%%%%
\section{Phase transition}
%%%
An important  consequence of  monotonicity and continuity of $\psi$ in $|\b|$ in Proposition \ref{prop:p} is the existence and uniqueness of %$\b_c^\pm$ defined below.
the critical temperatures introduced in the next statement, which is a direct consequence of the above.
%\begin{definition} 
%The critical (inverse) temperatures $\b_c^+(\nu), \b_c^-(\nu)$ at density $\nu$ are such that
%\begin{equation} \label{eq:highlow}
%\left\{ 
%\begin{array}{ccc}
%\psi(\b,\nu)=0  & {\rm if}  & \b \in [\b_c^-,\b_c^+]  \\
%\psi(\b,\nu)>0  & {\rm if}  & \b < \b_c^- \; {\rm or} \; \b > \b_c^+.  
%\end{array}
%\right.
%\end{equation}
%%We have $-\8 \leq \b_c^- \leq 0 \leq \b_c^+ \leq + \8$, and $|\b_c^\pm(\nu)|$ is non-increasing in $\nu$.
%%The domains in the $(\b,\nu)$-half-plane defined by the first and second line in \eqref{eq:highlow} are called {\rm high and low temperature region}
%%respectively.
%\end{definition}
\begin{theorem}\label{defdef..}
There exist $\b_c^+(\nu), \b_c^-(\nu)$ with  $-\8 \leq \b_c^- \leq 0 \leq \b_c^+ \leq + \8$ such that 
%
%The critical (inverse) temperatures $\b_c^+(\nu), \b_c^-(\nu)$ at density $\nu$ are such that
\begin{equation} \label{eq:highlow}
\left\{ 
\begin{array}{ccc}
\psi(\b,\nu)=0 \quad & {\rm if}  & \b \in [\b_c^-,\b_c^+]  \\
\psi(\b,\nu)>0 \quad  & {\rm if}  & \b < \b_c^- \; {\rm or} \; \b > \b_c^+.  
\end{array}
\right.
\end{equation}
%We have $-\8 \leq \b_c^- \leq 0 \leq \b_c^+ \leq + \8$, and $|\b_c^\pm(\nu)|$ is non-increasing in $\nu$.
%The domains in the $(\b,\nu)$-half-plane defined by the first and second line in \eqref{eq:highlow} are called {\rm high and low temperature region}
%respectively.
Moreover,  $|\b_c^\pm(\nu)|$ is non-increasing in $\nu$.
\end{theorem}
%We have $-\8 \leq \b_c^- \leq 0 \leq \b_c^+ \leq + \8$, and $|\b_c^\pm(\nu)|$ is non-increasing in $\nu$.
These values $\b_c^+(\nu), \b_c^-(\nu)$ are called \emph{ critical} (inverse)  \emph{ temperatures} at density $\nu$ (they depend on $r$ as well).
The domains in the $(\b,\nu)$-half-plane defined by the first and second line in \eqref{eq:highlow} are called {\em high and low temperature region}
respectively. The boundary between the two regions is called the {\em critical line}, and a {\em phase transition} in the statistical mechanics sense occurs:
the quenched free energy $p(\b,\nu)$ is equal to the annealed free energy $p^{(1)}(\b,\nu)=\nu \lambda r^d$ 
-- an analytic function -- but analyticity of  $p(\b,\nu)$ breaks down when crossing the critical line.  
\medskip


To summarize our finding, we define the  high temperature region and the  low temperature region
\begin{equation} \nn
{\cal D} = \{ (\b,\nu): \psi(\b, \nu)=0 \}, \qquad {\cal L}  = \{ (\b,\nu): \psi(\b, \nu)>0\}.
\end{equation}
They are are delimited by the critical lines $\b_c^-(\nu)$ and $\b_c^+(\nu)$ from Definition \ref{defdef..}.
In the next sections we will discuss non-triviality of the critical lines, as well as fine properties. In section \ref{ch_overlap} we will understand that they correspond to delocalized  or localized
behavior respectively.





%%%%%%%%%%%%%%%%%%%%%%
% We start with an informal description of  a specific example. 
%Consider a hydrophilic polymer chain (i.e., a long chain of monomers)  wafting in water. 
%Due to the thermal fluctuation, the shape of the polymer should be 
%understood as a random object. 
%We now suppose that the water contains randomly placed hydrophobic 
%molecules as impurities, which repel the 
%hydrophilic monomers which the polymer consists of. 
%The question we address here is:
%\begin{equation*} 
%\mbox{How does the impurities affect the global shape of the polymer chain?}
%\end{equation*}
%We try to answer this question in a mathematical framework.  
%However, as is everywhere else in mathematical physics, 
%it is very difficult to do so without compromising 
%with a rather simplified picture of the initial problem. 
%Here, our simplification goes as follows. 
%We first suppress entanglement, self-intersections  
%and U-turns of the polymer; we shall represent the polymer 
%chain as a graph $\{(j,x_j)\}^n_{j=1}$ in $\N \times \zd $, 
%so that the polymer is supposed to live in 
%$(1+d)$-dimensional discrete lattice 
%and to stretch in the direction of the first coordinate. Such a model is 
%called {\it directed}.
%Each point $(j,x_j) \in \N \times \zd $ on the graph 
%stands for the position of $j$-th monomer in this picture. Secondly, 
%we assume that, the transversal motion 
%$\bx= \{x_j\}^n_{j=1}$ performs a simple random walk in $\zd$, if the 
%impurities are absent. This accounts for consecutive monomers in the chain taking all possible configurations at a fixed distance one from another. 
%We then define the energy of the 
%path $\{(j,x_j)\}^n_{j=1}$ by 
%the formula (\ref{Ham}) at some (inverse) temperature 
% $\b>0$, where
% $\{\omega (n,x) : n \geq 1,\; x \in \zd  \}$ is a field of real 
%i.i.d.~random variables, with $\omega (n,x)$ describing the presence (or 
%strength) of an impurity  
%at site $(n,x)$. The case of an emulsion announced in the title is well captured by an i.i.d.~model. 
%Every time the path steps on an impurity it gets a penalty, and every step on a water molecule brings 
%a reward.
%The typical shape $\{(j,x_j)\}^n_{j=1}$ of the polymer 
%is then given by the one that maximizes the energy (\ref{Ham}).
% Let us suppose for example that 
%$\omega (n,x)$ takes two different values $+1$ (``presence of a water molecule 
%at $(n,x)$'') and $-1$ (``presence of the hydrophobic impurity 
%at $(n,x)$''). 
%Then, the energy of the polymer is decreased by 
%$\b$ each time a monomer is in contact with the impurity ($\omega (j,x_j)=-1$).
%Therefore, the typical shape of the polymer for each given 
%configuration of $\{ \omega (j, x)\}$ is given by the one 
%which tries to avoid the impurities as much as possible. 
%%%%%%%%%%
%\medskip
%
%{\bf What are the main questions ?} 
%We want to study large systems, that is, asymptotics when the length of the polymer chain goes to infinity,  for typical realizations of $\om$.  
%\medskip
%
%%%
%
%Mechanically, the elasticity of the polymer competes with the energy of very favorable, but distant, pinning sites that would require a costly distortion of the polymer.
%From this competition we can guess two possible scenarios:
%
%
%\begin{itemize}
%\item[]
%(i) If the space dimension is large 
%and the temperature is high, the impurities 
%should not affect the global shape of the polymer. This is because the polymer has much space 
%to avoid the impurities by the first assumption and the impurities are weak enough by the second one. It has no interest in travellling far, it only wants to pick a specific proportion (a little bigger than the density) of rewards, using a local strategy.
%\item[]
%(ii)
%If the dimension is small or the environment is strong, then the polymer will not be able any more to avoid the impurities if it stays at reasonable distances. The path will have an advantage in travelling far in order to find more favourable areas. 
%These atypical areas in the medium correspond to a higher density of rewards, and a precise geometry making them feasible for the walk. Then, the global shape of the polymer path 
%changes in a drastic manner. 
%\end{itemize}
%%%%%%%%%%%%%%%%%%%%%%%%%%%%%%%%%%%%%%%%%
%The informal description given above has been put into a mathematical framework by the 
%formalism of a Gibbs measure, in the framework of statistical mechanics. For a more substancial overview of polymer models, see the first chapter of 
%\cite{dHollanderStFl}. \index{modelization} \index{polymer model}
%%%%%%%%%
%\subsection{Polymers viewed as percolation at a positive temperature} \index{last passage  percolation} \index{percolation}
%\index{zero temperature} \index{ground states}
%
%For a fixed $n$, we can take the zero-temperature limits: as $\b \to \8$, 
%\begin{equation} \label{eq:0temp} \left\{
%\begin{array}{ccc}
%\b^{-1} \ln  Z_{n}(\om,\b) &  \longrightarrow &   \max_\bx  H_n(\bx) , \\
%\gibbs (\cdot)   & {\rm concentrates \ on}   & \; \arg \max_\bx   H_n(\bx) .
%\end{array}
%\right.
%\end{equation}
%The right-hand sides of \eqref{eq:0temp} are respectively the passage time %ground state energy
%and the set of geodesics of a last passage  percolation problem. Such quantities have  been extensively studied \cite{KestenStFlour, Durrett84, Newmantopics},
%and the last section of \cite{AuffDamronHanson50} for a recent account. Percolation here is oriented, site percolation.
%Thus, the polymer measure is an 
%interpolation between simple random walk (corresponding to $\b=0$) and  geodesics in last passage percolation (limit as $\b \to \8$). In particular it 
%can be viewed as a "fuzzy" geodesics, and it allows to study percolation by decreasing continuously 
%the temperature, performing a perturbative approach starting
%from the (well-known) random walk model. Having the parameter $\b$ at hand will reveal most useful. \index{perturbative}
%
%%%%%%%%%%%%%%%%%%%%%%%%% Fin CSYre
%
%\subsection{Exponents and localization} \index{characteristic exponents} \index{volume exponent} \index{wandering exponent} \index{roughness exponent}
%
%To quantify how the fluctuations are influenced by the disorder, two {\bf characteristic exponents}, 
%$\chi^\parallel \in [0,1/2]$ and $  \chi^\perp \in [1/2,1],$ are introduced: loosely speaking, 
%they are defined by
%\begin{equation} \label{def:exponents}
%P_n^{\b,\om} \left[ |S_n|^2 \right] \approx n^{2 \chi^\perp}\;,\quad {\rm Var}\big(\ln Z_n(\om,\b)\big) \approx n^{2\chi^\parallel } \;.
%\end{equation}
%The volume exponent $\chi^\parallel$ characterizes the fluctuations of thermodynamic quantities, though 
%the wandering exponent (or roughness exponent) characterizes the transverse fluctuations of the path. 
%It is expected, and partially confirmed (see \cite{AuffingerDamron13} adapting ideas  for last passage percolation summarized in \cite{Chatterjee-scaling}) 
%that the exponents are related via
%\begin{equation} \label{eq:exponentsrelated}
%\chi^\parallel = 2 \chi^\perp -1 \;.
%\end{equation}
%The scenario (i) in Section \ref{subsec:emulsion} above corresponds to diffusive behavior $\chi^\perp=1/2$, though case (ii)  
%has $\chi^\perp > 1/2$ (superdiffusivity). \index{superdiffusivity}
%In scenario (ii) these exponents are expected to be universal, i.e., not depending on the details of the model but only on the dimension, with 
%$\chi^\parallel =1/3,  \chi^\perp=2/3$ for $d=1$. 
%\medskip
%
%Another feature of scenario (ii) is {\bf localization}: the polymer has a large probability to lie in a favourable
%regions with high density of rewards and feasible for the path: we call them corridors, and this phenomenon is called path localization. It will be sudied in Chapters~\ref{ch:semimart} in a general framework and~\ref{ch:loggamma} on a specific integrable model, where explicit computations can be performed.
%\medskip
%
%
%The interest of mathematicians for these characteristic exponents, their relation and some inequalities depending on the space dimension, started two decades ago for last passage percolation in the discrete setting~\cite{NewPiza, LiceaNewPiza, Piza} and in the continuous setting~\cite{Wuthrich1, Wuthrich2}; it remains very strong
%with novel contributions \cite{AuffingerDamron-sc, Chatterjee-scaling}. For polymers, see also \cite{Petermann, Mejane, CaHu04, CYkokyuroku, BezerraTindelViens} for various estimates in the general setup. But except for exactly solvable models in $d=1$, they remain most mysterious - particularly for $d \geq 2$ - in spite of many efforts.
%\medskip
%
%
%
%%\medskip
%%
%%{\bf Variations}:
%%
%%The random walk $S$ can be defined on locally general finite graphs. In Section \ref{ch:tree},
%%it will be the simple random walk on a rooted binary tree, or more generaly, on a rooted regular tree. Such models relate to branching processes in inhomogeneous medium. 
%%Similarly, the model can be defined in continuous time and/or space.
%
%\section{Particle in a random potential}
%%
%We relate the Gibbs measure to branching random walks in an inhomogeneous medium.
%
%\subsection{Particle in deadly obstacles} \index{particle in deadly obstacles} \index{branching random walk} \index{Gibbs measure}
%
%\index{obstacles} \index{killing potential}
%In order to give a simple intuition of what is a Gibbs measure, 
%we first discuss here our model in the case when the environment variables
%are uniformly bounded from above. For simplicity, and without loss of generality\footnote{A constant shift for $\omega (t,x)$ has a trivial effect on the objects we consider.},
%we assume in this section that 
%\begin{equation}
%  \label{eq:hypneg}
%  \omega (t,x) \leq 0 \qquad \P-{\rm a.s.}
%\end{equation}
%In this case, we can interpret the model in terms of a walker moving among
%deadly obstacles. The space location $x$ hosts a deadly obstacle which strength
%at time $t$ is $(-\beta \omega (t,x))$. We start with an informal description. A particle starts at the origin at time 0,
%it jumps at integer times and immediately after, it dies or survives.
%It moves according to a simple random walk when alive, but it has a probability
%$\exp \beta \omega (t,x) \in (0,1]$ to survive the obstacle at time $t$ if it is still alive
%at time $t-1$ and has jumped at location $x$ at time $t$, and complementary probability
%$1-\exp \beta \omega (t,x)$ to be killed by the obstacle. The event of dying or 
%surviving at time $t$
%are independent of the past, conditionnally on being alive at time $t-1$
%and on the step at time $t$.
%\medskip
%
%We now give a precise construction of
%this process. Let $P, \P$ as before and let Unif be the uniform law
%on $[0,1]$. We fix a realization $(\omega (t,x); (t,x) \in \N\times \Z^d)$
%from $\P$. On the product space $\Omega_{\rm traj}  \times [0,1]$, we define the product measure
%$Pr=P \otimes {\rm Unif}$, and denote by $U:(\bx,u) \mapsto u$ the second
%coordinate mapping. Let
%$$
%\tau = \inf \big\{ n \geq 1: U > \exp\{\beta H_n^\om(S)\}\big\}
%$$
%be the lifetime of the path $S$. We also add to $\Z^d$ a special state called cemetery point and denoted by $\dagger$. 
%Then, the process defined by
%$$
%X_n= 
%\left\{
%  \begin{array}{ll}
%    S_n & {\rm if} \; n < \tau,\\
%    \dagger & {\rm if} \; n \geq \tau,
%  \end{array} \right.  
%$$
%corresponds to our informal description, and is intimately related to the polymer path.
%\begin{proposition}
%The sequence  $X$ is a (time-inhomogeneous) Markov chain on $\Z^d \cup \{ \dagger\}$. It has $\dagger$ as absorbing state, and its one-step transition probability from $x \in \Z^d$ is given by
%$$
%Pr(  X_{n+1}=y \vert X_{n}=x ) = 
%  (2d)^{-1}  \exp\{\beta \omega (n+1, y)\}\qquad {\rm if}\;  y \in \Z^d,
%$$
%and $$Pr(  X_{n+1}=\dagger \vert X_{n}=x ) =   (2d)^{-1}  \sum_{z \in \Z^d, |z-x|_1=1}
%\big(1-\exp\{\beta \omega (n+1, z)\}\big).$$
%Finally, the partition function of the polymer model  is equal to the survival probability,
%$$
%Pr( \tau > n ) = Z_n (\om,\beta)  ,
%$$
%and the polymer measure with time horizon $n$ is equal to the law 
%of the path of the particle killed in the random potential conditionned to be alive at time $n$,
%\begin{equation} \label{eq:avena2}
%\gibbs ( \bx)=
%Pr(X_{[0,n]}=\bx |  \tau > n ) 
%\;.
%\end{equation}
%\end{proposition}
%%%%%%
%%\begin{figure}[h]
%%\begin{center}
%%\includegraphics[scale=.15]{particle-in-potential.eps}
%%%%%\label{fig:ub} 
%%\caption{A particle in a deadly potential: The state space is $\Z^d \cup
%%\{\dagger\}$, with $d=2$ and $\dagger$ is a cemetery (absorbing) point. Conditionally on being still alive at time $t$ and at location 
%%$x$, it dies out there with probability $q=1-\exp \beta \omega (t,x)$ (recall 
%%that $\omega  \leq 0$), or survives one more time unit at least and 
%%moves to a neigboring site with probability $p=(2d)^{-1}(1-q)$. This description corresponds to a slightly 
%%different potential $\tilde H_n(S)= \sum_{i=0}^{n-1} \omega (i, S_i)$.
%%}
%%\end{center}
%%\end{figure}
%%%%%%%%% 
%\begin{proof}      
%First, we note that, for all $n$ and all path $\bx$ of length $n$,    
%\begin{eqnarray*}
%Pr( \tau > n | S_{[0,n]}=\bx) &=& {\rm Unif}(U \leq \exp\{\beta H_n(\bx)\})\\
%&=&
%\exp\{\beta H_n(\bx)\},
%\end{eqnarray*}
%and that
%\begin{equation}\label{eq:avena}
%Pr( \tau > n,  S_{[0,n]}=\bx) = (2d)^{-n}
%\exp\{\beta H_n(\bx)\},
%\end{equation}•
%and also that
%\begin{eqnarray*}
%Pr( \tau > n ) &=& Pr \big[ Pr( \tau > n | S_{[0,n]})\big]\\&=&
%Z_n (\om,\beta) .
%\end{eqnarray*}
%Therefore, 
%\begin{eqnarray*}
%Pr(X_{[0,n]}=\bx |  \tau > n ) &=& Pr(S_{[0,n]}=\bx |  \tau > n ) \\
%&=&
%\frac{Pr( \tau > n | S_{[0,n]}=\bx) Pr(S_{[0,n]}=\bx)}{Pr(\tau > n )}\\
%&=&
%\gibbs ( \bx),
%\end{eqnarray*}
%which is (\ref{eq:avena2}).
%We finally check that the dynamics of the random sequence $X=(X_n, n\geq 0)$ corresponds to the description at the beginning of the section.  First, $X$ is a Markov chain on $\Z^d \cup \{ \dagger\}$, with $\dagger$ as absorbing state.
%Moreover, for all nearest-neighbor path $\om$ in $\Z^d$ with $\bx_n=x, \bx_{n+1}=y$,
%\begin{eqnarray*}
%Pr\left( \tau \!>\! n\!+\!1,  S_{[0,n\!+\!1]}\!=\!\bx_{[0,n\!+\!1]} \!\!\!\!\!\!\!\!\! \right. &\vert&  \left.\!\!\!\!\!\!\!\!\!\tau > n, S_{[0,n]}\!=\!\bx_{[0,n]} \right)
%\\
%&=&\frac{Pr( \tau > n+1, S_{[0,n+1]}=\bx_{[0,n+1]}) }{Pr(\tau \!>\! n, S_{[0,n]}=\bx_{[0,n]} ) }\\
%&\stackrel{(\ref{eq:avena})}{=}& (2d)^{-1} \exp\{\beta (H_{n+1}(\bx)- H_n(\bx))\}
%\\&=&
% (2d)^{-1}  \exp\{\beta \omega (n+1, y)\}.
%\end{eqnarray*}
%This implies the first two equalities of the proposition. \qed \end{proof}
%%%%%%%%%%%%%%%%%%%%%%%%%%%%%%%
%%Except in the degenerate case $\b \omega (t,x)=0$ a.s., the lifetime $\tau$ is a.s. finite. Hence, the asymptotics of
%%the sequence of Gibbs measures $\gibbs $ does not simply reduce to $Pr( X \in \cdot | A)$, because the relevant 
%%$A=\{\tau<\8\}$ has probability zero here.
%%
%%\begin{remark}
%%When the environment variables are uniformly bounded from below.
%%\end{remark}
%
%%
%\subsection{Branching random walk picture} \label{sec:BRWp}
%
%We discuss now the case when the environment takes only nonnengative values.  We can interpret the model in terms of a branching population (without deaths) moving in a
%random potential:
%% 
%% The space location $x$ hosts a resource which allows a mean growth by a factor $\exp\{\beta \omega (t,x)\}$ of the local population. Here is an informal description of the dynamics of this random population.
%%We start with one individual at site $0$ at time 0, which will move and multiply from then on.
%%Each individual in the population jumps independently according to a simple random walk at integer times and immediately after, it branches (i.e., it is replaced by a random number of 
%%children) with mean value given by the above factor $\exp\{\beta \omega (t,x)\}$.
%%
%%Precisely, we first fix a collection $\nu_r$ indexed by $r \geq 1$ 
%%of probability measures on $\N^*$ with mean
%%$\sum_{k\geq 1} k \nu_r(k)=r$, such that the mapping $r \mapsto \nu_r$ is measurable.
%%An individual present in the population at time $t-1$ and location $y$ is replaced at time $t$  by 
%%a random number $k$ at a neighboring site $x \sim y$ with probability  
%%$(2d)^{-1} \nu_{\exp \b \omega (t,x)}(k)$.
%%
%%Then,  the partition function $Z_{\n,\beta}^\omega $ of the polymer model  is equal to the 
%% expected population size at time $n$ in a fixed environment $\omega $. 
%%The polymer measure with time horizon $n$ is equal to the law 
%%of the ancestral path of an individual selected uniformly at random in the full population at time $t$;
%The polymer measure with time horizon $n$ is equal to the law 
%of the ancestral path of an individual selected uniformly at random in the full population at time $t$.
%Moreover, the partition function $Z_{n,\beta}^\omega $ of the polymer model  is equal to the 
%expected population size at time $n$ in a fixed environment $\omega $. 
%
%%by ancestral path we mean the time-reversed of the path obtained by recording the location of the selected individual, the location of her father, 
%%the location of her grand-father, \dots till the location $0$ of her ancestor. 
%%\medskip
%%
%%We finally discuss  the general case without assumptions on the environment. We combine the above two constructions: fix a collection $\nu_r$ indexed by $r \geq 0$ 
%%of probability measures on $\N$ with mean
%%$\sum_{k\geq 0} k \nu_r(k)=r$, and construct a branching population moving like the simple random walk and branching according to these reproduction laws. When the outcome $k$ of the branching step is equal to 0, the individual dies without descendents. The whole population can get extinct.
%% The partition function $Z_{\n,\beta}^\omega $ of the polymer model  is still equal to the 
%% expected population size at time $n$ in a fixed environment $\omega $. 
%%But now, the polymer measure is equal to the law 
%%of the ancestral path of an randomly selected individual  at time $t$ conditionnally on survival
%%of the full population.
%\medskip
%
%The case of general environment is a simple combination of the two above: one can define a branching random walk in random environment such that, conditionally on survival at time $n$,  
%the law 
%of the ancestral path of an individual selected uniformly at random in the full population at that time
%is the polymer measure~\cite{CY11}.
%\medskip
%
%All these representations are in the spirit of Feynman-Kac formula.
%
%%%%%%%%%%%%%%%%%%%%%%%%%%
%
%\section{History, experiments and related models} \index{experiments}
%%%%%%%%%%
%The polymer model was originally introduced in physics literature \cite{HuHe85}
%to mimic the phase boundary of Ising model subject to random impurities. 
%It became a popular model for kinetic roughening of growing surfaces including front propagation, bacterial colonies, vortex lines in supraconductors,\ldots \cite{KrugSpohn}. Also, it has
%inspired a number of methods for its analysis (replica trick, renormalization group theory, Bethe ansatz, % [Brunet]
% scaling theory). 
%The model soon reached the mathematics community \cite{ImbrieSpencer88,Bolth89}, where it was explicitely formulated as the polymer measure  \eqref{mnen}.
%\medskip
%
%Let us mention  some physical realizations and real experiments fitting the model.
%
%\begin{itemize}
%\item[] (i) Torn paper sheets \cite{Kertesz}.
%A rectangular sheet of paper is tightened by a machine. The strain  is
%applied on two opposite sides, and it is slowly increased. A small notch 
%is made on a third side to initiate the tear. The fracture is governed 
%by the random geometry of the fiber network.   The fracture
%line is highly correlated to the weakest bonds in the sheet, as can be
%checked by microdensitometry. Directed polymers with $d=1$ and its zero
%temperature counterpart, oriented first
%passage percolation, are natural model for the fracture
%line. \\
%
%\item[] (ii) Nematic crystal: Takeuchi and Sano \cite{TakeuchiSano} realized an real experiment of a growing domain of topological-defect turbulence in electrically driven liquid crystal. 
%This gives an experimental evidence for KPZ scaling and Tracy-Widom fluctuations.
%Two videos of a flat or circular interface are available, see \\
%
%{\tt https://www.youtube.com/watch?v=tcdufma7UJA
%
%https://www.youtube.com/watch?v=useN8ZdRvb8}
%\end{itemize}
%
%The polymer model has a continuous limit expressed by a stochastic partial differential equation, the so-called Kardar-Parisi-Zhang
%equation (\ref{eq:KPZ}), which is the generic model for a variety of random growth phenomena and covers the statistics of roughness.
%The corresponding universality class consists in models whcih large-scale behavior share the same statistics.
%In spite of being 30 years old, KPZ equation remains of a tremendous interest in physics, as shown in recent  accounts~\cite{Halpin-healy} and  \cite{Thiery-these}.
%\medskip
%
%As mentioned above, corresponding to $\beta=\8$,  the ground states -- i.e. the paths $\bx$ 
%maximizing $H_n$ -- are the geodesics of last passage percolation.
%The relations of polymers to last passage percolation, extends to a number of related models:
%tandem queueing systems, totally asymmetric exclusion process, edge of the spectrum of random matrices, \ldots
%\medskip
%
%The parabolic Anderson model is a time-continuous version of our 
%discrete model, it is a popular example for intermittency phenomenon \cite{CarmonaMolchanov}. 
%It has been studied in many facets: fine asymptotics for vanishing diffusivity  \cite{CarmonaKoralovMolchanov, CMS02}, large deviations showing the asymmetry \cite{CranstonGM},
%using Gaussian tools \cite{CadelTindeViens}.  A complete review is \cite{GartnerKonig} in the time-independent case. An intermediate case is
%the so-called random walk in a dynamic random environment
%\cite{ErharddHollanderMaillard}
%%%
%\medskip
%
%To complete the section, we indicate a few models which are different but closely related: 
%%%
%\begin{enumerate}
%\item Polymer pinning: in  which localization only occurs at a known, specific location if it does occur.
%Main example: defect lines in the plane where rewards are available (and no rewards elsewhere). 
%Rewards may be deterministic or random.
%See \cite{Giacomin}, \cite{Giacomin2}, \cite{dHollanderStFl} for surveys of huge literature.
%\item Here, the polymer is given by a Brownian motion, whereas the environment is given by a space-time Poisson point process which is enlarged in the trasnsverse directions by a \. 
%All our forthcoming results have a counterpart here  \cite{CY05, CYkokyuroku, Va04}.  The model considered here enables to use stochastic calculus, with respect to both  Brownian motion and Poisson process, leading to handy formulas for fluctuations analysis and qualitative properties of the phase diagram. 
%The relation to  some formulation of the Kardar-Parisi-Zhang equation is transparent.
%Girsanov formula and the invariance of the Poisson process yields an explicit rate funtion for  large deviation of the path  \cite{CY05} and integration by parts yield complete localization
%\cite{CYBMPO2}. 
%\item In a similar spirit, it is natural to model the path itself by an independent increment process.  L\'evy polymers in random environment are considered in 
%\cite{Cstable, MTT}. Due to long jumps, the critical dimension can be different from that of diffusive models. 
%\item Random walks in a {\it time-independent} potential: spectral methods~\cite{Sznitman, Sznitman-milieu} are powerful in this case.
%The connection with our setup is when considering crossings of such walks when the path is stretched to a ballistic behaviour (bridge)~\cite{Ioffe-review, Sznitman}.
%Backtracks of the path are penalized, and the model we consider in this book becomes a reasonable approximation. In fact it corresponds to forbid backtracks and even more, to have a systematic move at each step in a  fixed direction. 
%\item The model here fits the paradigm of disordered systems and random medium: 
%local inhomogeneities with large scale regularity~\cite{Bovier} are nicely approximated by a random medium. 
%\item As mentioned in Section \ref{sec:BRWp} there is a strong connection with branching random walks in random environment. For a better focus we will not develop these 
%models. We simply mention here that, analogous to the results we discuss here, 
%diffusive regime \cite{HeilNaYo-clt, Shiozawa-clt, Yoshida08} and
%localized  behavior  \cite{HuYo09, Shiozawa-loc} have been studied in details.
%\item Linear stochastic evolution
%is a class of stochastic growth models on the integer lattice $\Z^d$, including a time discretization of the "linear systems'' discussed in \cite{LiggettIPS}. 
%Examples of models in this class are the number of open paths in oriented percolation and the binary contact path process. The counterparts of the topics we study here can be found in 
%\cite{FukushimaYoshida-ls, NagahataYoshida-loc, Yoshida-loc}.
%\item The effect of the correlated noise has been studied under various aspects: space-correlation  in the discrete case \cite{Lacoin-cor} and in the continuous case
%\cite{BTV, KMountford},
%time-correlation in the continuous case \cite{ViensZhang}.
%\item Localization covers the absence of diffusion in random medium: inhomogeneities can pin random paths and cancel their natural diffusive behavior.  Anderson's original 1958 paper addressed the question of spin diffusion on the lattice and electrical conduction in a semiconductor impurity band. The phenomenon has recieved 
%a huge interest to understand its mechanism as well as for its applications in science and material ingeneering. We refer to the reviews
%\cite{Fifty, Thouless} for a physics account, and \cite{Hundermark} for a mathematical introduction to localization  in Anderson model.
%\end{enumerate}
%
%
%\section{Highlights} %of the following chapters}
%
%
%
%
%Accessible and building on general methods, the monograph is designed to be read linearly.
%The main topics are organized as described below.
%\medskip
%
%In Chapter 2 we take the view point of statistical mechanics for disordered systems. We introduce the free energy,  prove self-averaging, and determine the qualitative shape the phase diagram and find a thermodynamic phase transition on the basis of a comparison with an annealed model. As stated in Theorem \ref{beta_c}, the critical temperature $\b_c$ is determined by the onset of strict inequality between 
%the quenched and annealed rate of growth of the partition function,
%\begin{eqnarray}
%\b > \b_c & \iff& \lim_n n^{-1} \ln Z_n(\om, \b) <  \lim_n n^{-1} \ln \P Z_n(\om, \b) ,\\
%0 \leq \b \leq \b_c & \iff& \lim_n n^{-1} \ln Z_n(\om, \b) =  \lim_n n^{-1} \ln \P Z_n(\om, \b) .
%\end{eqnarray}
%In the Chapter 3 we take a more probabilistic approach, showing that martingale theory yields direct information on the rescaled partition function. For instance $W_n=Z_n(\om, \b)/\P[Z_n(\om, \b)]$ is a positive, mean-one martingale.
%The well-known "second moment method" reveals efficient in a large regime, where analytic  methods can be used. 
%Computing integer moments, we will check that the martingale $W_n$ is bounded in $L^2$ when $d$ is large and 
%$\b$ is small, implying that the phase transition occurs in a non-trivial manner in dimension 3 and above (cf. Theorem \ref{th:exists-WD}). Sharp results 
%can be obtained in this region, such as the existence of a finite random  limits for 
%$\gibbs[|S_n|^2]-n$ or $\gibbs[H_n]-n \lambda'(\b)$.
%
%
%
%In Chapter 4 we introduce the polymer models on trees - or, equivalently, the branching random walk in random environment,   which can be seen as a mean-field approximation.
%This model is   solvable, and we will discuss similarities and  differences with the polymer on the lattice. 
%
%In Chapter 5 we use semi-martingale techniques
%for a fine analysis of the weak and the strong disorder regimes, relating these regimes to the fine behavor of the random polymer. Standard tools from stochastic processes, as quadratic variations, yield quantitative information on
%the overlap between replicas, i.e., independent copies of the polymer path sharing the same environment.
%This is the key to understanding that the thermodynamic phase transition coincide with the localization transition.   
%We will obtain in Theorem \ref{loc-deloc} that $\b > \b_c$ is equivalent to the positivity of the Cesaro limit of the probability of the favorite end-point of the polymer,
%$$
%J_n = \max\{ \gibbs (S_n=x); x \in \Z^d\} . 
%$$
%Thus, the critical temperature is the onset of localization of the end-point of the polymer.
%
%In Chapter 6 we study the localized phase. 
%It happens at all temperature in dimension 1 and 2. We describe the corridors where the polymer concentrates.
%The last section will deal with the simpler situation where the medium has heavy tails. 
%Then, a coarse graining is performed to analyze localization.
%
%
%In Chapter 7 we study in details the log-gamma polymer recently introduced by Sepp\"al\"ainen, which is integrable 
% in one dimension.
%Explicit computations are possible, and we will review some of the precise results from an abundant literature,  including characteristic exponents and  fluctuation statistics. Slightly modifying the model on the boundaries we will have the benefit of  stationarity, and we will completely determine the limiting distribution of $J_n$, cf. Corollary \ref{cor:first}.
%Also, in this model with boundaries, we will obtain tightness of the polymer 
%around the favorite endpoint, and the large scale fluctuation of the favorite endpoint. 
%The mechanism for localizing the polymer is by confinement in the deepest valley a random potential - here, a random walk. We believe that this mechanism, which appears in Sinai's localization for the  random walk in random environment in the one-dimensional recurrent case~\cite{Sinai82}, 
%is an important ingredient of 
%polymer localization in general.
%
%Chapter 8 will be dedicated to the KPZ equation in dimension $d=1$.  We will discuss what is the  meaning of the equation,  express its solution as a series, define the continuum random polymer that it generates, and study its relation to intermediate disorder for polymers. 
%We will show convergence of the rescaled, discrete polymer in dimension $d=1$ and with a  suitably
%decaying inverse temperature $\b=\b_n$, to the  continuum random polymer.
%
%In Chapter 9 we will discuss the variational formulas giving the value of free energy. 
%They come in two forms: optimizing over 
%cocycles as in stochastic homogeneization, or optimizing over probability measures
%as in usual Gibbs variational formulas, giving a balance between energy and entropy. 
%
%Finally, an appendix collects a few general tools - mainly from stochastic processes and statistical mechanics - which
%are used  all through the monograph.
%
